




\section{Related Work}\label{sec:related}
\subsection{Coding Agents \& Code as Policies}
\label{sec:related-coding}
This line of work rests on a central abstraction---\emph{executable code as the agent's action}: the model generates a program, the runtime returns structured feedback, and the loop iterates. Code-as-Policies \citep{cap} and ProgPrompt \citep{singh2022progprompt} introduced this formulation for robotics, composing perception and skill APIs into a task plan, and visual reasoning carried the same compositional idea to vision modules \citep{suris2023vipergpt}. Subsequent research extended single-shot generation into multi-turn loops driven by execution feedback---reason-and-act traces \citep{yao2022react}, self-critique and iterative repair \citep{shinn2023reflexion, madaan2023self}, inference-time search \citep{yao2023tree}, and learned tool invocation \citep{schick2023toolformer}; \citet{wang2024codeact} argues code is a stronger action representation than language because each step is runtime-verifiable. A parallel thread formalizes the \emph{agent--computer interface}, a sandboxed shell in which agents read, run, and debug code, now underlying frontier coding products \citep{openai2026codex, anthropic2026claudeopus47, moonshot2026kimicode}.

Within robotics, early code-as-policy systems exposed high-level, human-engineered APIs---closed-form skills, language-conditioned grasping primitives, and affordance scorers \citep{saycan}---leaving the agent to decompose tasks rather than synthesize low-level control. \citet{wang2023voyager} grows a reusable Minecraft skill library via self-verification and an automatic curriculum where trials are near-free, and \citet{capx2026} finds that multi-turn feedback, skill synthesis, and ensembled sampling improve manipulation reliability over low-level primitives. A related line uses LLMs to synthesize auxiliary training signals: reward functions \citep{ma2024eureka, xie2024text2reward, yu2023language}, sim-to-real transfer protocols \citep{ma2024dreureka}, simulation environments \citep{wang2023robogen}, and data-collection pipelines \citep{ahn2024autort}.

%Two assumptions underlie this body of work: that trials are inexpensive, and that the agent's code controls the policy alone. Our setting violates both. Real-robot time is the binding budget, and the agent's code drives the entire research loop---constructing the experimental environment, developing reusable perception and manipulation tools, and scheduling training experiments---rather than only the policy. The relevant constraint accordingly shifts from the expressiveness of the action space, which has been the focus of prior code-as-policy work, to the cost of each physical execution.


\subsection{Agentic Self-Improvement}\label{sec:related-selfimprove}
A self-improvement loop needs each trial to be cheap enough to repeat at scale; systems differ mainly in what a trial returns to the agent. \citet{ellis2021dreamcoder} set the retention pattern that the rest build on: each solved synthesis task adds named subroutines to a growing library, which is viable because execution is free. \citet{wang2023voyager} carried this into an embodied setting, swapping the success signal for LLM self-verification and an automatic curriculum---again viable because Minecraft rollouts cost nothing. \citet{yu2023language}, \citet{ma2024eureka}, and \citet{xie2024text2reward} return a reward rather than a skill: an LLM proposes a dense reward, a policy is trained against it, and the reward is revised from training statistics---a loop Eureka closes via thousands of Isaac Gym rollouts per minute. \citet{ma2024dreureka} reaches toward hardware via synthesized domain-randomization parameters, yet still iterates in simulation and deploys only after revision stops; at the experimental-setting level, \citet{wang2023robogen} generates new tasks and assets in simulation, and \citet{ahn2024autort} coordinates a mobile-robot fleet for offline data collection rather than within-loop iteration. In every case, the loop closes on a cheap substrate while real-robot execution serves only as a sim-to-real or data target, never as the medium of iteration. We retain these skill-accumulation and reward-generation mechanisms but run the loop directly on hardware, where the binding resource is the agent's robot-access budget, not its compute.


\subsection{Autonomous Research Agents and Scientific Discovery}
\label{sec:related-research} The final body of work automates the research loop itself, best read along \emph{the medium in which experiments are launched}. Prior to LLMs, autonomous systems closed the hypothesis--experiment loop on real laboratory hardware \citep{king2004functional, burger2020mobile}. LLM-era systems instead automate the \emph{digital} loop end to end \citep{lu2024aiscientist, yamada2025aiscientistv2, schmidgall2025agentlab, schmidgall2025agentrxiv}, with a recent strand reaching back into the physical sciences through lab automation \citep{boiko2023autonomous, m2024augmenting} or human-executed wet-lab validation \citep{ghareeb2025robin}. Evaluation tracks this digital emphasis: MLE-bench \citep{chan2025mle} scores ML-engineering agents and resource scaling, while analyses of SWE-bench \citep{jimenez2024swebench} caution that gains can reflect contamination over capability. Two gaps follow: no system autonomously runs and optimizes a physical \emph{robotics} loop under an explicit hardware budget---classical robot scientists fixed the apparatus and never wrote their own tools, while LLM research agents never touch real robots---and no benchmark measures utilization of a scarce physical resource rather than capability or cost-per-paper. Our system addresses both gaps by running and optimizing the experimental loop on real robots and introducing resource-utilization metrics for this budget-bound regime.




% Coding agents share a common operating principle: the model emits code, observes execution feedback, and revises. One thread establishes the loop as a general reasoning-and-acting mechanism and contributes its now-standard ingredients: self-critique, iterative refinement, candidate search, and tool invocation \citep{yao2022react, shinn2023reflexion, madaan2023self, yao2023tree, schick2023toolformer}. In the digital domain, this loop can be built as an explicit \emph{agent--computer interface} ---a sandboxed terminal in which agents read, run, and debug code \citep{jimenez2024swebench, yang2024swe, wang2025openhands,
% wang2024codeact}. Recent works consider extending coding agents to the physical domain, specializing it to robotics via \emph{code-as-policy}, where the agent's actions are executable program that composes perception and control primitives \citep{liang2023cap,
% singh2022progprompt, suris2023vipergpt}. CaP-X \citep{capx2026} unifies the first two for manipulation and shows that the third's lesson---scaling test-time interaction with visual differencing, skill synthesis, and ensembling---recovers reliability even over low-level primitives. What unites these threads is an environment where execution feedback is effectively free. 
%Our harness inherits the loop but closes it over physical hardware, where the cost of each trial, not the expressiveness of the action space, is the binding constraint.

